\documentclass[11pt]{article}
\usepackage[margin=1in]{geometry}
\usepackage[T1]{fontenc}
\usepackage[utf8]{inputenc}
\usepackage{lmodern}
\usepackage{amsmath,amssymb,amsthm,mathtools}
\usepackage{booktabs,array,longtable}
\usepackage{enumitem}
\usepackage{setspace}
\usepackage{microtype}
\usepackage{csquotes}
\usepackage{hyperref}
\usepackage[backend=biber,style=authoryear,natbib=true,maxcitenames=2,maxbibnames=99]{biblatex}
\usepackage{titlesec}
\hypersetup{colorlinks=true,linkcolor=blue,citecolor=blue,urlcolor=blue}
\setstretch{1.08}
\newtheorem{proposition}{Proposition}
\newtheorem{lemma}{Lemma}
\newtheorem{assumption}{Assumption}
\newcommand{\E}{\mathbb{E}}
\newcommand{\R}{\mathbb{R}}
\newcommand{\Prob}{\mathbb{P}}
\newcommand{\one}{\mathbf{1}}
\DeclareMathOperator*{\argmax}{arg\,max}
\titleformat{\section}{\large\bfseries}{\thesection}{0.7em}{}
\titleformat{\subsection}{\normalsize\bfseries}{\thesubsection}{0.7em}{}
\titleformat{\subsubsection}{\normalsize\itshape}{\thesubsubsection}{0.7em}{}

\addbibresource{Sanctions_Risk_Dollar_Convenience.bib}

\title{Sanctions Risk and Dollar Convenience: A Top-level PhD Research Proposal}
\author{Prepared as an independent proposal draft}
\date{July 2026}

\begin{document}
\maketitle
\begin{abstract}
This standalone proposal (Original Topic 2) states a focused research question, positions it against verified literature, develops a tractable model, derives testable implications, and specifies an empirical design to validate or reject the mechanism. It is self-contained: the preamble is embedded in this file and the references are contained in the local BibTeX file only.
\end{abstract}

\paragraph{Source-integrity note.} References are restricted to publications, working papers, datasets, and institutional sources checked during preparation. Unverified claims are not used as established facts.

\section{Sanctions Risk and the Endogenous Erosion of Dollar Convenience}

\paragraph{One-sentence pitch.} Financial sanctions are powerful because the dollar system is dominant, but every use of sanctions may reduce the perceived safety of dollar assets. This project builds a dynamic reserve-portfolio model in which sanctions risk endogenously erodes the convenience yield of the dominant currency.

\subsection*{1. Research question and reviewer positioning}
The question is precise:
\begin{quote}
How does the anticipation of future financial sanctions change sovereign reserve portfolios, the dollar convenience yield, and the sanctioning country's optimal use of financial infrastructure?
\end{quote}
This is not a broad paper on de-dollarization. It is a paper about a specific asset-pricing wedge: the value of dollar liquidity net of state-contingent expropriation or freezing risk.

The factual motivation is clear. The freezing of Russian central-bank reserves after the 2022 invasion of Ukraine made the safety of reserve assets explicitly state-contingent. Brookings reports that the REPO Task Force estimated frozen reserves at about $280$ billion, while many experts place the amount between $300$ and $330$ billion \citep{BrookingsRussiaAssets2025}. Recent theory already establishes the relevance of this mechanism. \citet{BianchiSosaPadilla2024JME} model restrictions on international reserves as sanctions and show how reserve freezes interact with sovereign default. \citet{BianchiSosaPadilla2025EJ} show that anticipated financial sanctions can reduce the dollar convenience yield and dollar-asset holdings. \citet{GopinathStein2021} provide the complementary dominant-currency mechanism: trade invoicing and safe-asset demand reinforce one another.

The proposed contribution is to endogenize the dynamic feedback from the sanctioning country's policy to the future size of the dollar network. The sanctioning country chooses a sanction intensity today; reserve holders update perceived freezing risk and adjust portfolios; portfolio adjustment lowers the convenience yield generated by dollar network externalities; lower convenience yield reduces the future power of sanctions.

\subsection*{2. Contribution relative to the literature}
The paper extends three literatures. First, it builds on the sovereign-sanctions models of \citet{BianchiSosaPadilla2024JME} and \citet{BianchiSosaPadilla2025EJ}. Second, it embeds sanctions risk in the dominant-currency logic of \citet{GopinathStein2021}. Third, it uses sanctions data such as the Global Sanctions Data Base of \citet{FelbermayrEtAl2020} and reserve data from IMF COFER \citep{IMFCOFER}. The contribution is not to document that reserves are diversified. The contribution is to connect reserve diversification to an equilibrium policy tradeoff for the hegemon.

A top-field referee will demand a clean margin. The clean margin is: a dollar asset pays a liquidity service flow but carries a state-contingent accessibility risk for countries exposed to sanctions. That wedge can be priced, estimated, and disciplined.

\subsection*{3. Model}
There is a continuum of countries $i\in[0,1]$ and a dominant-currency issuer $H$. Country $i$ chooses the dollar share $a_{i,t}\in[0,1]$ of reserves. Non-dollar reserves are normalized as an outside asset. Let
\begin{equation}
N_t=\int_0^1 a_{i,t}\,di
\end{equation}
be the size of the dollar network. A dollar reserve asset delivers financial return $r^D_t$ and liquidity services $\ell(N_t)$, with $\ell'>0$ and $\ell''<0$. The outside reserve asset returns $r^O_t$ and liquidity service normalized to zero. Country $i$ faces sanctions probability
\begin{equation}
\pi_{i,t}=\pi(z_{i,t},g_{i,t};S_t),
\end{equation}
where $z_{i,t}$ is its geopolitical exposure, $g_{i,t}$ is its action or alignment relative to $H$, and $S_t$ is the sanctioning country's policy rule. If sanctioned, fraction $\varphi\in[0,1]$ of dollar reserves becomes inaccessible for the relevant state.

Country $i$ solves
\begin{equation}
\max_{a_{i,t}\in[0,1]} \; a_{i,t}\left(r_t^D-r_t^O+\ell(N_t)-\pi_{i,t}\varphi\right)-\frac{\kappa_i}{2}a_{i,t}^2-\frac{\nu}{2}(a_{i,t}-a_{i,t-1})^2. \label{eq:reserve-choice}
\end{equation}
The quadratic terms capture diversification costs and adjustment costs. The interior first-order condition is
\begin{equation}
a_{i,t}=\frac{r_t^D-r_t^O+\ell(N_t)-\pi_{i,t}\varphi+\nu a_{i,t-1}}{\kappa_i+\nu}. \label{eq:reserve-FOC}
\end{equation}
Equation \eqref{eq:reserve-FOC} is the empirical object: dollar demand falls with sanctions risk and rises with the endogenous convenience yield.

The hegemon receives per-period payoff
\begin{equation}
U^H_t=\underbrace{\Omega(N_t)}_{\text{exorbitant privilege}}+\underbrace{G(s_t,g_{T,t})}_{\text{geopolitical benefit}}-\underbrace{L(s_t,D_t)}_{\text{default and market cost}}-\underbrace{C(N_t,N_{t+1})}_{\text{network erosion}}, \label{eq:hegemon}
\end{equation}
where $s_t\in[0,1]$ is sanction intensity against target $T$, and $D_t$ is the target's default indicator or debt-service disruption. The law of motion for the network is
\begin{equation}
N_{t+1}=\int_0^1 A_i\left(r_t^D-r_t^O+\ell(N_t)-\pi(z_{i,t},g_{i,t};s_t)\varphi,a_{i,t}\right)di. \label{eq:network-law}
\end{equation}

\subsection*{4. Main theoretical results to prove}
\begin{proposition}[Sanctions-risk wedge]
For any country with interior reserve share, an increase in expected sanctions loss $\pi_{i,t}\varphi$ lowers the dollar reserve share one-for-one scaled by $(\kappa_i+\nu)^{-1}$, holding the network convenience yield fixed.
\end{proposition}

\begin{proposition}[Strategic complementarity in dollar holding]
If $\ell'(N)>0$, then reserve choices are strategic complements: a fall in aggregate dollar holdings lowers the dollar convenience yield and induces additional reserve diversification by other countries.
\end{proposition}

\begin{proposition}[Time inconsistency of sanctions]
Suppose $G_s>0$, $G_{ss}<0$, and sanctions reduce $N_{t+1}$ through \eqref{eq:network-law}. The ex post optimal sanction intensity exceeds the ex ante intensity that maximizes the discounted value of the dollar network whenever the marginal geopolitical benefit is high and the network erosion cost is convex.
\end{proposition}

The third result is the central policy insight. A dominant-currency issuer has an option to weaponize financial infrastructure. Exercising the option yields immediate geopolitical payoff but reduces the future value of the option by damaging the network.

\subsection*{5. Empirical design}
\paragraph{Data.} The empirical strategy combines IMF COFER aggregate reserve composition \citep{IMFCOFER}, country-level US Treasury holdings from TIC data, sanctions episodes from \citet{FelbermayrEtAl2020}, public measures of frozen Russian sovereign assets \citep{BrookingsRussiaAssets2025}, UN voting alignment from \citet{BaileyStrezhnevVoeten2017}, and controls for exchange-rate valuation effects, trade exposure, commodity prices, and global risk.

\paragraph{Design 1: reserve portfolio response to financial-sanction exposure.} The baseline panel is
\begin{equation}
\Delta DollarShare_{i,t}=\alpha_i+\tau_t+\beta\,Exposure_i\times PostSanctionShock_t+\Gamma'X_{i,t}+\varepsilon_{i,t}. \label{eq:sanction-did}
\end{equation}
$Exposure_i$ can be measured by geopolitical distance from the sanctioning coalition, historical sanctions exposure, similarity to the sanctioned target in foreign-policy alignment, or military/diplomatic ties. $PostSanctionShock_t$ is defined around major financial-sanctions events, especially reserve freezes and payment restrictions. The coefficient $\beta<0$ is predicted by the model.

\paragraph{Design 2: convenience yield channel.} Estimate whether changes in sanctions salience move proxies for the dollar convenience yield. Possible objects include Treasury convenience-yield measures, deviations from covered interest parity, and relative yields on safe assets. A reduced-form specification is
\begin{equation}
CY_t=\alpha+\rho CY_{t-1}+\delta SanctionSalience_t+\Xi'Z_t+u_t, \label{eq:convenience}
\end{equation}
where $SanctionSalience_t$ is an index constructed from financial-sanctions announcements and media intensity. The model predicts $\delta<0$ for salience that raises perceived accessibility risk.

\paragraph{Design 3: structural estimation.} The parameters $(\varphi,\kappa_i,\nu,\ell'(N))$ are estimated by matching reserve-share responses, adjustment persistence, and convenience-yield movements. The structural model then runs counterfactuals:
\begin{enumerate}[label=(\roman*)]
\item no reserve freeze in 2022;
\item soft restrictions only;
\item repeated sanctions against multiple targets;
\item a larger non-dollar outside asset market.
\end{enumerate}

\subsection*{6. Linking empirics to the model}
The project must validate the theory on three distinct margins. First, exposed countries should reduce dollar holdings more after financial-sanctions shocks. Second, the response should be stronger when the sanctioned asset is closer to a reserve asset, not merely when ordinary trade sanctions occur. Third, reserve diversification by exposed countries should have spillover effects on non-exposed countries only through the convenience-yield channel. These three margins separate the proposed mechanism from generic de-dollarization narratives.

\subsection*{7. Falsifiable predictions}
\begin{enumerate}[label=\textbf{P\arabic*.},leftmargin=*]
\item Countries more exposed to future sanctions risk reduce dollar reserve shares or US Treasury holdings after salient financial-sanctions events.
\item The effect is stronger for countries with greater geopolitical distance from the sanctioning coalition.
\item Ordinary trade sanctions should have weaker reserve-portfolio effects than reserve freezes or payment-system restrictions.
\item Dollar convenience-yield proxies decline when financial-sanctions salience rises, controlling for global risk and monetary policy.
\item The sanctioning country's optimal policy is less aggressive when the network elasticity $\ell'(N)$ is high.
\end{enumerate}

\subsection*{8. Dissertation structure}
\begin{enumerate}[leftmargin=*]
\item \textbf{Chapter 1: Theory.} Build the dynamic reserve-portfolio and hegemon-sanction model; prove strategic complementarity and time inconsistency.
\item \textbf{Chapter 2: Empirics.} Estimate reserve-portfolio responses to financial-sanctions shocks using exposed versus non-exposed countries.
\item \textbf{Chapter 3: Counterfactual policy.} Quantify the optimal sanctions rule under network erosion and evaluate welfare effects for the hegemon, targets, and third countries.
\end{enumerate}

\subsection*{9. Reviewer stress test}
The main danger is overreach. The paper should not claim to explain the entire future of the international monetary system. It should claim to identify and quantify one wedge: sanctions-adjusted liquidity value. The second danger is identification, because reserve shares change slowly and many reserve managers do not disclose country-level currency composition. The fallback is to use TIC holdings, gold accumulation, and valuation-adjusted aggregate COFER moments as complementary evidence. The third danger is conflating sanctions risk with geopolitical risk generally. The empirical design must distinguish reserve freezes and payment restrictions from standard trade sanctions.

\subsection*{10. Bottom line}
This proposal has high frontier value but needs strict narrowing. A top version is not ``financial sanctions and de-dollarization.'' It is ``sanctions risk as a state-contingent haircut on reserve liquidity.'' That is a precise economic object and a credible dissertation chapter.


\newpage
\printbibliography
\end{document}
